\section{模 Abel 簇}
\label{sec:modular-abelian-varieties}
% ============================================================================

到这里，本章的桥已经搭好了一半：
$J_0(N)$ 把整块空间 $\Sk_2(\GammaO(N))$ 装进一个 Abel 簇里，Hecke 代数 $\mathbb T$ 也已经进入 $\End(J_0(N))$。
但我们真正关心的通常不是整块 Jacobian，而是其中与\emph{某一个}新形式 $f$ 对应的那部分几何信息。

\textbf{我们遇到了什么问题？}
$J_0(N)$ 往往维数很大，包含了许多不同特征形式混合在一起的信息。
如果我们只想研究某个归一化特征新形式 $f$，就必须找到一种办法，把“所有 Hecke 特征值等于 $a_n(f)$ 的部分”从整个 Jacobian 里切出来。
否则，模形式与椭圆曲线之间的对应仍然会被其他特征形式干扰。

\textbf{这个概念的核心思想是什么？}
Hecke 代数已经作用在 $J_0(N)$ 上，因此我们可以把那些在 $f$ 上作用为零的 Hecke 元全部收集成一个理想
$I_f$，再把 $J_0(N)$ 对这个理想生成的子簇取商。
这个商恰好保留了 $f$ 的特征值系统，而杀掉了所有不属于 $f$ 的方向。

\textbf{引入之后能得到什么？}
这样构造出的商 Abel 簇 $A_f$ 就是\textbf{模 Abel 簇}。
它的维数等于 Fourier 系数域 $K_f=\Q(a_n(f):n\ge 1)$ 的次数。
特别地，当 $K_f=\Q$ 时，$A_f$ 维数为 $1$，所以它是一条椭圆曲线。
这就是从单个新形式直接提取出椭圆曲线的机制。

\input{figures/ch06-modular-abelian}

\begin{definition}[特征值理想]
\label{def:ideal-If}
设 $f\in \Sk_2(\GammaO(N))$ 是归一化 Hecke 特征新形式，
并设
\[
  K_f=\Q\bigl(a_n(f):n\ge 1\bigr)
\]
为其 Fourier 系数生成的数域。
由\cref{prop:eigenform-gives-character}，存在环同态
\[
  \lambda_f:\mathbb T\longrightarrow K_f,
  \qquad T_n\longmapsto a_n(f).
\]
定义
\[
  I_f=\ker \lambda_f.
\]
它称为与 $f$ 相关的\textbf{Hecke 理想（Hecke ideal）}\index{Hecke理想!Hecke ideal}。
\end{definition}

\begin{definition}[模 Abel 簇]
\label{def:modular-abelian-variety}
设 $J_0(N)$ 上由 $I_f$ 生成的 Abel 子簇记为 $I_fJ_0(N)$。
定义
\[
  A_f=J_0(N)/I_fJ_0(N).
\]
它称为与 $f$ 对应的\textbf{模 Abel 簇（modular abelian variety）}\index{模Abel簇!modular abelian variety}。
\end{definition}

\begin{proposition}[模 Abel 簇的余切空间]
\label{prop:cotangent-of-Af}
设 $q_f:J_0(N)\twoheadrightarrow A_f$ 是商映射。
则拉回映射给出单射
\[
  q_f^*:H^0(A_f,\Omega^1)\hookrightarrow H^0(J_0(N),\Omega^1)
  \cong \Sk_2(\GammaO(N)),
\]
其像恰为
\[
  \Sk_2(\GammaO(N))[I_f]
  =\{h\in \Sk_2(\GammaO(N)):\ Th=\lambda_f(T)h\text{ 对所有 }T\in\mathbb T\}.
\]
也就是被理想 $I_f$ 消掉的 Hecke 特征子空间。
\end{proposition}

\begin{proof}
商映射 $q_f$ 是满的全纯群同态，所以拉回 $q_f^*$ 在微分上必为单射。

下面确定其像。
设 $\omega\in H^0(A_f,\Omega^1)$，则对任意 $T\in I_f$，因为 $T$ 在商 $A_f$ 上作用为零，
有
\[
  T^*(q_f^*\omega)=q_f^*(T^*\omega)=0.
\]
把 $H^0(J_0(N),\Omega^1)$ 与 $\Sk_2(\GammaO(N))$ 识别后，这说明 $q_f^*\omega$ 被每个 $T\in I_f$ 消掉，
故其像包含在 $\Sk_2(\GammaO(N))[I_f]$ 中。

反过来，设 $h\in \Sk_2(\GammaO(N))$ 被所有 $T\in I_f$ 消掉。
在 Lie 代数层面，$A_f$ 的切空间是 $T_0J_0(N)$ 模去 $I_f$ 生成的子空间；
对偶地，余切空间正是被 $I_f$ 消掉的线性泛函构成的子空间。
因此这样的 $h$ 正好来自商的余切空间，也就是属于 $q_f^*$ 的像。
故像恰为所述子空间。

最后，若 $h$ 同时被所有 $T-\lambda_f(T)$ 消掉，就等价于
\[
  Th=\lambda_f(T)h \qquad (T\in\mathbb T),
\]
这正是右边的描述。
\end{proof}

\begin{theorem}[模 Abel 簇的维数]
\label{thm:dim-Af}
设 $f\in \Sk_2(\GammaO(N))$ 是归一化 Hecke 特征新形式，系数域为 $K_f$。
则
\[
  \dim A_f=[K_f:\Q].
\]
更精确地说，$H^0(A_f,\Omega^1)$ 由 $f$ 的所有 Galois 共轭
\[
  f^\sigma \qquad (\sigma:K_f\hookrightarrow \C)
\]
张成。
\end{theorem}

\begin{proof}
由\cref{prop:cotangent-of-Af}，
$H^0(A_f,\Omega^1)$ 嵌入到 $\Sk_2(\GammaO(N))$ 中，且恰好等于被理想 $I_f$ 消掉的子空间。

先证明每个 Galois 共轭 $f^\sigma$ 都在这个子空间里。
对任意 $T\in\mathbb T$，若 $\lambda_f(T)\in K_f$，则
\[
  T(f^\sigma)=\sigma\bigl(Tf\bigr)=\sigma\bigl(\lambda_f(T)f\bigr)=\sigma(\lambda_f(T))f^\sigma.
\]
特别地，当 $T\in I_f$ 时 $\lambda_f(T)=0$，故 $Tf^\sigma=0$。
因此所有 $f^\sigma$ 都属于 $\Sk_2(\GammaO(N))[I_f]$。

再证明这里没有别的特征方向。
由第\ref{ch:hecke}章的自伴性与可交换性，$\Sk_2(\GammaO(N))$ 可分解为 Hecke 特征形式的直和。
设 $g$ 是其中一个特征形式，且被 $I_f$ 消掉。
若 $g$ 不是 $f$ 的 Galois 共轭，则存在某个 Hecke 算子 $T_n$ 使得
\[
  a_n(g)\neq \sigma\bigl(a_n(f)\bigr)
\]
对所有嵌入 $\sigma$ 都成立；特别地 $a_n(g)\neq a_n(f)$。
于是元素
\[
  T_n-a_n(f)\in I_f
\]
在 $g$ 上的作用为
\[
  (T_n-a_n(f))g=(a_n(g)-a_n(f))g\neq 0,
\]
与“$I_f$ 消掉 $g$”矛盾。
所以被 $I_f$ 消掉的特征方向只能来自 $f$ 的 Galois 共轭。

因此
\[
  H^0(A_f,\Omega^1)=\operatorname{span}_\C\{f^\sigma\}.
\]
这些共轭彼此线性无关：若存在关系
\[
  c_1f^{\sigma_1}+\cdots+c_rf^{\sigma_r}=0,
\]
比较某个 $a_n(f)$ 的各共轭作用，便得到一个 Vandermonde 型线性系统，只能使所有 $c_i=0$。
更直接地说，不同嵌入作用在系数域上得到不同的 Fourier 系数列，故对应的 $q$-展开不同，从而线性无关。

嵌入 $\sigma:K_f\hookrightarrow\C$ 的个数正是 $[K_f:\Q]$，
所以
\[
  \dim H^0(A_f,\Omega^1)=[K_f:\Q].
\]
Abel 簇的维数等于其全纯微分空间维数，因此
\[
  \dim A_f=[K_f:\Q].
\]
\end{proof}

\begin{corollary}[有理系数情形给出椭圆曲线]
\label{cor:kf-q-elliptic}
若 $K_f=\Q$，则 $A_f$ 是一维 Abel 簇，因此是一条椭圆曲线。
\end{corollary}

\begin{proof}
由\cref{thm:dim-Af}，此时
\[
  \dim A_f=[K_f:\Q]=1.
\]
再由\cref{prop:dim1-torus-abelian}，一维 Abel 簇就是椭圆曲线。
\end{proof}

\begin{example}[级别 $11$ 的模椭圆曲线]
\label{ex:level11-modular-abelian}
在\cref{ex:jacobian-motivation-11,ex:j011-elliptic} 中我们已经知道
\[
  \dim \Sk_2(\GammaO(11))=1,
  \qquad
  \dim J_0(11)=1.
\]
因此唯一的归一化新形式
\[
  f=q-2q^2-q^3+2q^4+q^5+\cdots
\]
满足 $K_f=\Q$。
由\cref{cor:kf-q-elliptic}，对应的 $A_f$ 是椭圆曲线。
又因为 $J_0(11)$ 本身已经一维，商映射
\[
  J_0(11)\twoheadrightarrow A_f
\]
只能是同种，实际上是同构。
所以 $J_0(11)$ 与该新形式对应的模椭圆曲线是同一个对象。
经典模型可取为
\[
  E: y^2+y=x^3-x^2-10x-20.
\]
本章不证明这个 Weierstrass 方程的具体推出过程；
第\ref{ch:eichler-shimura-l}章将解释它与 $f$ 的 $L$-函数为何完全一致。
\end{example}

\textbf{Eichler--Shimura 预告。}
对权 $2$ 新形式 $f$，第\ref{ch:eichler-shimura-l}章会证明
\[
  L(A_f,s)=\prod_{\sigma:K_f\hookrightarrow\C} L(f^\sigma,s).
\]
特别地，当 $K_f=\Q$ 时，$L(A_f,s)=L(f,s)$。
这就是“特征形式 $\leftrightarrow$ 椭圆曲线”最精确的解析表达。
